\section{Consequences}

Let $V, W$ be representations of a Lie algebra $\g$.

\begin{claim}
$\g \to \End(V \otimes W) = \End(V) \otimes \End(W)$ via
$x \mapsto x \otimes 1 + 1 \otimes x$ is a homomorphism of Lie algebras.
\end{claim}

\begin{exercise}
Prove it. (This comes from differentiating the group homomorphism
$G \to G \times G$ via $g \mapsto (g, g)$.)
\end{exercise}

\begin{corollary}
If $V, W$ are representations of $\g$, then so is $V \otimes W$.
\end{corollary}

\begin{remark}
If $A$ is an algebra and $V$ and $W$ are representations of $A$, then
$V \otimes W$ is naturally a representation of $A \otimes A$. To make it a
representation of $A$, we need an algebra homomorphism $A \to A \otimes A$.
(An object with such a structure, plus some other properties, is called a Hopf
algebra.) In some essential way, $\g$ behaves like a Hopf algebra. (Or,
rather, the universal enveloping algebra of $\g$ is a Hopf algebra, and so are
its deformations, the quantum groups.)
\end{remark}

Now, take $\g = \sl_2$. What is $L(n) \otimes L(m)$? That is, we know
$L(n) \otimes L(m) = \bigoplus_k a_kL(k)$; what are the $a_k$?

One method of doing this would be to compute all the highest-weight vectors.

\begin{exercise}
Do this for $L(1) \otimes L(n)$. Then do it for $L(2) \otimes L(n)$.
\end{exercise}

As a start on the exercise, let $v_a$ denote the highest-weight vector in
$L(a)$. Then $v_n \otimes v_m$ is a highest-weight vector in
$L(n) \otimes L(m)$. Indeed,
\begin{align*}
  h(v_n \otimes v_m) &= (hv_n) \otimes v_m + v_n \otimes (hv_m) = (n + m)\,v_n \otimes v_m,\\
  e(v_n \otimes v_m) &= (ev_n) \otimes v_m + v_n \otimes (ev_m) = 0 + 0 = 0.
\end{align*}
Therefore,
\[
  L(n) \otimes L(m) = L(n + m) + \text{other stuff}.
\]
However, counting dimensions, we get
\[
  (n + 1)(m + 1) = (n + m + 1) + \dim(\text{other stuff}),
\]
which shows that there is quite a lot of this ``other stuff'' in there.

One can write down explicit formulae for all the highest-weight vectors; these
are complicated but mildly interesting. However, we don't have to do it to
determine the summands of $L(n) \otimes L(m)$.

Let $V$ be a finite-dimensional representation of $\sl_2$.

\begin{definition}
The \emph{character} of $V$ is
$\ch V = \sum_{n \in \Z} \dim V_n\, z^n \in \N[z, z^{-1}]$.
\end{definition}

\begin{properties}\leavevmode
\begin{enumerate}
  \item $\ch V|_{z=1} = \dim V$.

    This is equivalent to the claim that $h$ is diagonalisable with
    eigenvalues in $\Z$, so $V = \bigoplus_{n \in \Z} V_n$.
  \item $\ch L(n) = z^n + z^{n-2} + \dots + z^{-n}
       = \dfrac{z^{n+1} - z^{-(n+1)}}{z - z^{-1}}$.
    Sometimes this is written as $[n + 1]_z$.
  \item $\ch V = \ch W$ iff $V \cong W$.

    Note that $\ch L(0) = 1$, $\ch L(1) = z + z^{-1}$,
    $\ch L(2) = z^2 + 1 + z^{-2}, \dots$ form a $\Z$-basis for the space of
    symmetric Laurent polynomials with integer coefficients. On the other
    hand, by complete reducibility,
    \begin{align*}
      V &= \textstyle\bigoplus_{n \geq 0} a_nL(n), \quad a_n \geq 0,\\
      W &= \textstyle\bigoplus_{n \geq 0} b_nL(n), \quad b_n \geq 0,\\
      V \cong W &\iff a_n = b_n \text{ for all } n.
    \end{align*}
    But since the $\ch L(n)$ form a basis, $\ch V = \sum a_n \ch L(n)$
    determines the $a_n$.

    \begin{remark}
    Note that for the representations, we only need nonnegative coefficients
    $a_n \geq 0$. On the other hand, if instead of looking at modules we were
    looking at chain complexes, the idea of $\Z$-linear combinations would
    become meaningful.
    \end{remark}
  \item $\ch(V \otimes W) = \ch V \cdot \ch W$.

    \begin{exercise}[Essential!]
    $V_n \otimes W_m \subseteq (V \otimes W)_{n+m}$. Therefore,
    $(V \otimes W)_p = \bigoplus_{n+m=p} V_n \otimes W_m$. This is exactly how
    we multiply polynomials.
    \end{exercise}

    \begin{example}
    We can now compute $L(1) \otimes L(3)$:
    \begin{align*}
      \ch L(1) \cdot \ch L(3) &= (z + z^{-1})(z^3 + z + z^{-1} + z^{-3})\\
        &= (z^4 + z^2 + 1 + z^{-2} + z^{-4}) + (z^2 + 1 + z^{-2})
         = \ch L(4) + \ch L(2),
    \end{align*}
    so $L(1) \otimes L(3) = L(4) \oplus L(2)$.
    \end{example}
  \item The Clebsch--Gordan rule:
    \[
      L(n) \otimes L(m) = \bigoplus_{\substack{k = |n - m|\\ k \equiv n - m \bmod 2}}^{n + m} L(k).
    \]
    There is a picture that goes with this rule (so you don't actually need to
    remember it); see figure~\ref{fig:crystal-sl2}.
\end{enumerate}
\end{properties}

\begin{figure}[!htbp]
  \centering
  \slcrystal
  \caption{Crystal for $\sl_2$, drawn for $n = 4$, $m = 6$.}
  \label{fig:crystal-sl2}
\end{figure}

The individual strings are the $L(a)$ in the decomposition.

One of the points of this course is that this sort of picture-drawing (called
crystals) lets us decompose tensor products of representations of all other
semisimple Lie algebras. (This fact was discovered about 20 years ago by
Kashiwara and Lusztig.)

What we will actually show for the other semisimple Lie algebras:
\begin{itemize}
  \item The category of finite-dimensional representations is completely
    reducible (i.e.\ all finite-dimensional representations are semisimple).
  \item We will parametrise the irreducible representations.
  \item We will compute the characters of the irreducible representations.
  \item We will show how to take tensor products of irreducible
    representations by drawing pictures (crystals).
\end{itemize}
